\clearpage
\section{계산 절차와 장 요약}
\label{sec:norm-trace-workflow-review}

\begin{learninggoal}
노름·대각합 문제를 읽고 가장 짧은 계산 경로를 선택한다. 분리성, 갈루아성, 확장탑, 유한체 여부를 차례로 점검하는 재현 가능한 절차를 익힌다.
\end{learninggoal}

\subsection{문제 해결 절차}

\begin{keyidea}
노름과 대각합은 먼저 ``무엇을 알고 있는가''를 분류한 뒤 계산한다. 행렬을 무조건 만드는 것은 가장 확실하지만 항상 가장 빠른 방법은 아니다.
\end{keyidea}

\subsubsection*{1단계: 확장과 기준체를 명시한다}

반드시
\[
  \operatorname N_{L/K},
  \qquad
  \Tr_{L/K}
\]
에서 $L$과 $K$를 확인한다. 같은 원소라도 더 큰 체에서 계산하면 켤레가 반복되어 값이 달라진다. 예를 들어
\[
  \operatorname N_{\Q(\sqrt2)/\Q}(\sqrt2)=-2
\]
이지만
\[
  \operatorname N_{\Q(\sqrt2,\sqrt3)/\Q}(\sqrt2)=4
\]
이다.

\subsubsection*{2단계: 가장 가까운 구조를 찾는다}

\begin{enumerate}[label=(\alph*)]
  \item 기저와 곱셈표가 주어졌는가?
  
  $m_\alpha$의 행렬을 만들고 행렬식과 대각합을 취한다.

  \item 최소다항식이 주어졌는가?
  
  $P_{\alpha,L/K}=f_\alpha^{[L:K(\alpha)]}$를 사용한다.

  \item $g(\alpha)$의 노름인가?
  
  $L=K(\alpha)$이면
  \[
    \operatorname N(g(\alpha))=\Res(f_\alpha,g)
  \]
  를 먼저 시도한다.

  \item 확장이 분리 또는 갈루아인가?
  
  모든 매장 또는 갈루아 자동동형에 대한 합과 곱을 사용한다.

  \item 중간체가 있는가?
  
  추이식
  \[
    \Tr_{L/K}=\Tr_{M/K}\circ\Tr_{L/M},
  \]
  \[
    \operatorname N_{L/K}
    =\operatorname N_{M/K}\circ\operatorname N_{L/M}
  \]
  로 나눈다.

  \item 유한체인가?
  
  프로베니우스 공식
  \[
    \Tr(x)=\sum_{j=0}^{n-1}x^{q^j},
    \qquad
    \operatorname N(x)=x^{(q^n-1)/(q-1)}
  \]
  을 사용한다.
\end{enumerate}

\subsubsection*{3단계: 분리성을 확인한다}

분리확장이면 켤레들의 합과 곱을 그대로 사용할 수 있다. 비분리확장이면 서로 다른 매장만 더해서는 안 되고 비분리차수의 중복도를 포함해야 한다. 특히 유한 비분리확장에서는
\[
  \Tr_{L/K}=0
\]
이다.

\subsubsection*{4단계: 결과를 독립적으로 검산한다}

다음 검산을 사용할 수 있다.

\begin{enumerate}[label=(\roman*)]
  \item 노름은 곱셈적이어야 한다.
  \item 대각합은 $K$-선형이어야 한다.
  \item $c\in K$에 대해
  \[
    \operatorname N(c)=c^{[L:K]},
    \qquad
    \Tr(c)=[L:K]c
  \]
  여야 한다.
  \item 갈루아 자동동형을 적용해도 값이 변하지 않아야 한다.
  \item 확장탑으로 계산한 값과 직접 계산한 값이 같아야 한다.
  \item 최소다항식의 상수항과 둘째 계수의 부호를 확인한다.
\end{enumerate}

\subsection{대표 판정표}

\begin{center}
\renewcommand{\arraystretch}{1.28}
\begin{tabular}{p{0.28\textwidth}p{0.60\textwidth}}
\toprule
질문 & 사용할 사실 \\
\midrule
$\Tr(\alpha)$와 $\operatorname N(\alpha)$ & $m_\alpha$의 대각합과 행렬식 \\
최소다항식만 알고 있음 & $P_{\alpha,L/K}=f_\alpha^{[L:K(\alpha)]}$ \\
분리확장 & 모든 $K$-매장에 대한 합과 곱 \\
갈루아 확장 & $\Gal(L/K)$의 모든 원소에 대한 합과 곱 \\
확장탑 & 노름·대각합의 추이성 \\
비분리확장 & 비분리차수를 중복도로 포함; 대각합은 영 \\
대각합 쌍의 비퇴화성 & 유한확장의 분리성과 동치 \\
$g(\alpha)$의 노름 & $\Res(f_\alpha,g)$ \\
유한체 & 프로베니우스 합과 지수 \\
순환확장의 노름 $1$ & $a/\sigma(a)$, 힐베르트 정리 90 \\
순환확장의 대각합 $0$ & $a-\sigma(a)$, 가법형 힐베르트 정리 90 \\
\bottomrule
\end{tabular}
\end{center}

\subsection{자주 하는 실수}

\begin{pitfall}
\textbf{실수 1: 최소다항식의 차수와 전체 확장차수를 혼동한다.}

$\alpha$가 $L$ 전체를 생성하지 않으면 특성다항식은 최소다항식 자체가 아니라
\[
  f_\alpha^{[L:K(\alpha)]}
\]
이다. 노름에는 거듭제곱이, 대각합에는 배수가 생긴다.
\end{pitfall}

\begin{pitfall}
\textbf{실수 2: 노름의 부호를 빠뜨린다.}

일계수 최소다항식
\[
  T^n+a_{n-1}T^{n-1}+\cdots+a_0
\]
에 대해
\[
  \operatorname N(\alpha)=(-1)^n a_0,
  \qquad
  \Tr(\alpha)=-a_{n-1}.
\]
특히 홀수차에서는 상수항의 부호가 바뀐다.
\end{pitfall}

\begin{pitfall}
\textbf{실수 3: 비분리확장에서 서로 다른 매장만 곱하거나 더한다.}

서로 다른 매장의 개수는 분리차수뿐이다. 특성다항식에는 비분리차수만큼의 중복도가 추가된다.
\end{pitfall}

\begin{pitfall}
\textbf{실수 4: 대각합의 전사성과 $\Tr(1)$을 동일시한다.}

분리확장에서는 대각합이 전사이지만 $\Tr(1)$은 표수 때문에 $0$일 수 있다.
\end{pitfall}

\begin{pitfall}
\textbf{실수 5: 힐베르트 정리 90을 노름의 전사성으로 오해한다.}

정리는 노름이 $1$인 원소를 매개화한다. 일반 $c$가 노름인지 여부는 별도의 산술 문제이다.
\end{pitfall}

\subsection{개념 연결 지도}

\[
\begin{array}{c}
\text{체의 원소 }\alpha\\[2pt]
\downarrow\\[-2pt]
\text{곱셈선형사상 }m_\alpha
\end{array}
\qquad
\begin{array}{c}
\xrightarrow{\det}\operatorname N_{L/K}(\alpha)\\[10pt]
\xrightarrow{\operatorname{tr}}\Tr_{L/K}(\alpha)
\end{array}
\]

분리확장에서는
\[
  m_\alpha\text{의 고유값}
  =\sigma_1(\alpha),\dots,\sigma_n(\alpha),
\]
따라서
\[
  \det=\text{켤레들의 곱},
  \qquad
  \operatorname{tr}=\text{켤레들의 합}.
\]
체매장들의 선형독립은 대각합 쌍의 비퇴화성을 만들고, 이는 분리성을 판정한다. 순환확장에서는 다시 힐베르트 정리 90이 노름과 대각합의 핵을 설명한다.

\begin{chapterreview}
\begin{enumerate}[label=\arabic*.]
  \item $\operatorname N_{L/K}(\alpha)$와 $\Tr_{L/K}(\alpha)$는 곱셈선형사상 $m_\alpha$의 행렬식과 대각합이다.
  \item $P_{\alpha,L/K}=f_\alpha^{[L:K(\alpha)]}$이므로 최소다항식의 계수에서 노름과 대각합을 읽을 수 있다.
  \item 유한 분리확장에서는 노름과 대각합이 모든 체매장에 의한 켤레들의 곱과 합이다.
  \item 갈루아 확장에서는 그 합과 곱을 갈루아군 전체에 대해 취한다.
  \item 확장탑에서 노름과 대각합은 각각 합성된다.
  \item 비분리차수는 매장 공식에 공통 중복도로 나타나며, 유한 비분리확장의 대각합은 영사상이다.
  \item 서로 다른 캐릭터와 체매장은 함수로서 선형독립이다.
  \item 대각합 쌍 $(x,y)\mapsto\Tr(xy)$은 정확히 분리확장에서 비퇴화이다.
  \item 유한 분리확장의 대각합은 전사지만 노름은 일반적으로 전사가 아니다.
  \item 유한체에서는 두 사상이 모두 전사이고 명시적인 프로베니우스 공식과 핵 크기를 갖는다.
  \item $L=K(\alpha)$에서 $\operatorname N(g(\alpha))=\Res(f_\alpha,g)$이다.
  \item 순환확장에서 힐베르트 정리 90은 노름 $1$과 대각합 $0$인 원소를 각각 $a/\sigma(a)$와 $a-\sigma(a)$로 나타낸다.
\end{enumerate}
\end{chapterreview}

\begin{checkpoint}
\begin{enumerate}[label=(\alph*)]
  \item 최소다항식, 갈루아군, 확장탑이 모두 주어진 문제에서 어떤 계산법을 먼저 선택할지 기준을 세워라.
  \item 분리성 판정에 대각합 사상 자체가 아니라 대각합 쌍이 필요한 이유를 설명하라.
  \item 노름·대각합·판별식·종결식이 하나의 특성다항식 관점에서 어떻게 연결되는지 서술하라.
\end{enumerate}
\end{checkpoint}
